
\documentstyle[12pt]{article}
%%%%%%%%%%%%%%%%%%%%%%%%%%%%%%%%%%%%%%%%%%%%%%%%%%%%%%%%%%%%%%%%%%%%%%%%%%%%%%%%%%%%%%%%%%%%%%%%%%%%%%%%%%%%%%%%%%%%%%%%%%%%
%TCIDATA{OutputFilter=LATEX.DLL}
%TCIDATA{LastRevised=Wednesday, December 19, 2001 09:34:27}
%TCIDATA{<META NAME="GraphicsSave" CONTENT="32">}

    
\topmargin -1.7cm
    \oddsidemargin -0.05cm 
    \evensidemargin -0.05cm
    \textheight 23.6cm 
    \textwidth 16.5cm 
    \sloppy
    
\def\newline{\hfil\break}
    \font\tr=ptmr at 12pt 
    \font\trb=ptmb at 12pt
    
\input{tcilatex}

\begin{document}


\thispagestyle{empty} \tr %\draft
\noindent

\begin{center}
{\sc Generalized Spin Density Functional Theory: Yamaguchi Papers}\\[0pt]
\vspace{0.3cm}  S.B.~Trickey \\[0pt]
\vspace{0.3cm}  18 Dec. 2001 \\[0pt]
\vspace{0.6cm}
\end{center}

In early Fall 2001 my attention was drawn to two references in a paper  that
suggested that Yamaguchi had introduced generalized spin orbitals  and
Fukutome classification in DFT as far back as 1979, substantially  before
Weiner and Trickey, Internat J. Quantum Chem. {\bf 69}, 451-60 (1998).  The
references given were  \vspace{-0.2cm}

\begin{itemize}
\item K. Yamaguchi, Chem. Phys. Lett. {\bf 67}, 75 (1979)

\item {\em ibid} {\bf 69}, 111 (1979).
\end{itemize}

B.~Weiner made an effort to get these references but found that the  volume
and page numbers did not correspond to any article by Yamaguchi,  with or
without co-authors. On Dec. 7, 2001, I undertook a laborious  search in the
library, with the following results:

\begin{enumerate}
\item Weiner is correct; neither of the aforementioned citations 
corresponds to papers.

\item No citation to Yamaguchi appears in either  vol. 67 or 69 of {\em %
Chemical Physics Letters}.

\item After extensive searching that is, however, not guaranteed to  be
exhaustive, I was able to find and read the following:  \vspace{-0.2cm}

\begin{enumerate}
\item K. Yamaguchi, T. Fueno, and H. Fukutome, Chem. Phys. Lett. {\bf 22}, 
461-65 (1973)

\item K. Yamaguchi, Chem. Phys. Lett. {\bf 33}, 330-35 (1975)

\item K. Yamaguchi, Chem. Phys. Lett. {\bf 35}, 230-35 (1975)

\item K. Yamaguchi and T. Fueno, Chem. Phys. Lett. {\bf 38}, 47-51 (1976)

\item K. Yamaguchi and T. Fueno, Chem. Phys. Lett. {\bf 38}, 52-56 (1976)

\item K. Yamaguchi and T. Fueno, Chem. Phys. {\bf 19}, 35-42 (1977)

\item K. Yamaguchi and T. Fueno, Chem. Phys. {\bf 29}, 117-39 (1978)

\item K. Yamaguchi, Chem. Phys. Lett. {\bf 66}, 395-99 (1979)

\item K. Yamaguchi, Chem. Phys. Lett. {\bf 68}, 477-82 (1979)
\end{enumerate}

\item Of these, all but the last two are about UHF, GSOs, etc.,  and have no
direct bearing on DFT (except for comparisons). Only  the last two have
anything to do with DFT and then {\em only}  in the context of so-called
Hartree-Fock-Slater (HFS) theory. Paper "i"  is focused exclusively on the
use of the HFS determinant as  a reference (and the orbitals as a basis) for
CI, coupled-cluster, and  MBPT calculations. Therefore, only Paper "h" is
relevant to the  history of GSOs and Fukutome classification in DFT.
Presumably  "h" is the paper that the incorrect references I had encountered
intended to cite. Why there are two incorrect references yet  only one paper
is unclear.

\item Paper "h" did {\em not} establish a GSO formulation of Kohn-Sham
theory. Rather, it  proceeded by simple analogy to Hartree-Fock, under the
long-since-abandoned  assumption that HFS is approximate HF. Note that the 
single determinant wave functions play distinctly different roles in  HF
theory and in DFT. Therefore, Paper "h" provided {\em no} rationale  related
to DFT for introducing the GSOs nor did it treat the  implications for the
relationship between the exchange-correlation  energy functional and the
consequent KS equation. Therefore "h"  is not the original formulation for a
GSO treatment of KS theory,  though it is admittedly a precursor. The major
steps taken  by Weiner and Trickey  were to show that the logic of the KS
scheme {\em requires} the use of  GSOs and the Fukutome classification if
full generality of spin  configurations is to be achieved and, therefore,
that the  exchange-correlation kernels {\em must} be derived specifically to
each Fukutome class.
\end{enumerate}

\end{document}
